% !TEX root = ../../ctfp-print.tex

\lettrine[lhang=0.17]{想}{为一本讲范畴论的书收尾}, 并没有一个好地方. 总有更多的东西要学. 范畴论是
一门浩瀚的学问. 与此同时, 同样的一些主题、概念和模式显而易见地一遍又一遍地冒出来. 有一种说法是:
一切概念都是 Kan 扩张 —— 事实上, 你确实可以用 Kan 扩张推出极限、余极限、伴随、单子、米田引理, 以及
多得多的东西. 范畴这个概念本身在各个抽象层次上都会出现, 幺半群和单子的概念也是如此. 哪一个才是
最根本的呢? 结果发现, 它们全都相互关联, 一个引出另一个, 构成一个永无止境的抽象循环. 我决定,
用展示这些相互关联来为这本书收尾, 也许是个好办法.

\section{双范畴}

范畴论最困难的方面之一, 就是视角的不断切换. 以集合范畴为例. 我们习惯于用元素来定义集合. 空集没有元素.
单元素集合有一个元素. 两个集合的笛卡尔积是一个序对构成的集合, 如此等等. 但谈到范畴 $\Set$ 的时候,
我请你忘掉集合的内容, 转而关注它们之间的态射 (箭头). 你可以时不时地掀开盖子看一眼, 看看 $\Set$ 里
某个特定的万有构造用元素说出来是什么样子. 终止物件原来就是含一个元素的集合, 如此等等. 但这些只是
些健全性检查.

函子被定义成范畴之间的一个映射. 把映射看作范畴里的态射是很自然的. 函子原来就是范畴范畴
(如果想回避大小的问题, 就说小范畴范畴) 里的一个态射. 一旦把函子当成箭头, 我们就放弃了关于它对范畴内
部(物件与态射) 的作用的信息, 就像我们在 $\Set$ 里把函数当成箭头时, 放弃了关于它对集合中元素的作用的
信息一样. 但任意两个范畴之间的函子本身也构成一个范畴. 这一次, 要你来把某个在一个范畴里是箭头的东西,
看成另一个范畴里的物件. 在函子范畴里, 函子是物件, 自然变换是态射. 我们发现, 同一个东西可以在
一个范畴里是箭头, 而在另一个范畴里是物件. 把物件当名词、把箭头当动词这种朴素的看法是站不住的.

与其在两种视角之间来回切换, 我们不如试着把它们合并成一个. 这样就得到了 $\cat{2}$-范畴的概念: 其中的
物件叫 $0$-胞腔, 态射叫 $1$-胞腔, 而态射之间的态射叫 $2$-胞腔.

\begin{figure}[H]
  \centering
  \includegraphics[width=0.35\textwidth]{images/twocat.png}
  \caption{$0$-胞腔 $a, b$; $1$-胞腔 $f, g$; 以及一个 $2$-胞腔 $\alpha$.}
\end{figure}

\noindent
范畴范畴 $\Cat$ 就是一个现成的例子. 我们有作为 $0$-胞腔的范畴, 作为 $1$-胞腔的函子, 以及作为
$2$-胞腔的自然变换. $\cat{2}$-范畴的定律告诉我们: 任意两个 $0$-胞腔之间的 $1$-胞腔构成一个范畴
(换句话说, $\cat{C}(a, b)$ 是一个 hom-范畴, 而不是 hom-集合). 这跟我们早先那个论断 ——
任意两个范畴之间的函子构成一个函子范畴 —— 正好吻合.

特别地, 从某个 $0$-胞腔回到它自己的 $1$-胞腔也构成一个范畴, 即 hom-范畴 $\cat{C}(a, a)$;
但这个范畴的结构还要更多. $\cat{C}(a, a)$ 的成员既可以被看作 $\cat{C}$ 里的箭头, 也可以被
看作$\cat{C}(a, a)$ 里的物件. 作为箭头, 它们可以彼此复合. 但当我们把它们看成物件时, 复合就变成了
一个从一对物件到一个物件的映射. 事实上, 它看起来非常像一个积 —— 准确地说, 是一个张量积.
这个张量积有一个单位: 恒等 $1$-胞腔. 结果发现, 在任何 $\cat{2}$-范畴里, hom-范畴 $\cat{C}(a, a)$
自动就是一个幺半范畴, 其张量积定义为 $1$-胞腔的复合. 结合律和单位律直接由相应的范畴定律落出来.

我们来看看这在 $\cat{2}$-范畴的典范例子 $\Cat$ 里意味着什么. hom-范畴 $\Cat(a, a)$ 就是 $a$ 上的
自函子范畴. 自函子的复合在其中扮演张量积的角色. 恒等函子是这个积的单位. 我们先前见过,
自函子构成一个幺半范畴 (我们在定义单子时用过这个事实), 但现在我们看到这是一个更普遍的现象:
任何$\cat{2}$-范畴里的自 $1$-胞腔都构成一个幺半范畴. 等我们推广单子的时候, 还会回到这一点.

你也许还记得, 在一般的幺半范畴里, 我们并不要求幺半群定律严格成立. 单位律和结合律往往在同构的
意义下成立就够了. 在 $\cat{2}$-范畴里, $\cat{C}(a, a)$ 里的幺半群定律来自 $1$-胞腔的复合定律.
这些定律是严格的, 所以我们总会得到一个严格幺半范畴. 不过, 这些定律也是可以放松的. 比如我们可以说:
恒等 $1$-胞腔 $\idarrow[a]$ 与另一个 $1$-胞腔 $f \Colon a \to b$ 的复合, 同构于 (而不是等于) $f$.
$1$-胞腔的同构是用 $2$-胞腔定义的. 换句话说, 存在一个 $2$-胞腔:
\[\rho \Colon f \circ \idarrow[a] \to f\]
它有一个逆.

\begin{figure}[H]
  \centering
  \includegraphics[width=0.35\textwidth]{images/bicat.png}
  \caption{双范畴里的恒等律在同构的意义下成立 (一个可逆的
    $2$-胞腔 $\rho$).}
\end{figure}

\noindent
对左单位律和结合律, 我们可以如法炮制. 这种放松了的 $\cat{2}$-范畴叫作双范畴 (还有一些附加的相容性定律
,
这里我就不写了).

正如预料的那样, 双范畴里的自 $1$-胞腔构成一个带有非严格定律的一般幺半范畴.

双范畴的一个有趣的例子是 span 范畴. 两个物件 $a$ 与 $b$ 之间的一个 span 是一个物件 $x$ 和一对态射:
\begin{gather*}
  f \Colon x \to a \\
  g \Colon x \to b
\end{gather*}

\begin{figure}[H]
  \centering
  \includegraphics[width=0.35\textwidth]{images/span.png}
\end{figure}

\noindent
你也许还记得, 我们在定义范畴积时用过 span. 这里, 我们想把 span 看成双范畴里的 $1$-胞腔. 第一步是
定义span 的复合. 假设我们有一个相邻的 span:
\begin{gather*}
  f' \Colon y \to b \\
  g' \Colon y \to c
\end{gather*}

\begin{figure}[H]
  \centering
  \includegraphics[width=0.5\textwidth]{images/compspan.png}
\end{figure}

\noindent
这个复合会是第三个 span, 带某个顶点 $z$. 对它来说最自然的选择, 就是 $g$ 沿 $f'$ 的拉回. 记住,
一个拉回是一个物件 $z$ 连同两个态射:
\begin{align*}
  h  & \Colon z \to x \\
  h' & \Colon z \to y
\end{align*}
使得:
\[g \circ h = f' \circ h'\]
它在所有这类物件中是万有的.

\begin{figure}[H]
  \centering
  \includegraphics[width=0.5\textwidth]{images/pullspan.png}
\end{figure}

\noindent
眼下我们只关注集合范畴上的 span. 在这种情况下, 拉回就是从笛卡尔积 $x \times y$ 里取出的、满足下式的
那些序对 $(p, q)$ 构成的集合:
\[g\ p = f'\ q\]
共享同一对端点的两个 span 之间的态射, 定义为它们顶点之间的一个态射 $h$, 使得相应的三角形可交换.

\begin{figure}[H]
  \centering
  \includegraphics[width=0.4\textwidth]{images/morphspan.png}
  \caption{$\cat{Span}$ 里的一个 $2$-胞腔.}
\end{figure}

\noindent
总结一下, 在双范畴 $\cat{Span}$ 里: $0$-胞腔是集合, $1$-胞腔是 span, $2$-胞腔是 span 态射. 恒等
$1$-胞腔是一个退化的 span, 其中三个物件全都相同, 而那两个态射都是恒等态射.

我们以前还见过双范畴的另一个例子: \hyperref[ends-and-coends]{前函子}构成的双范畴 $\cat{Prof}$, 其中
0-胞腔是范畴, 1-胞腔是前函子, 2-胞腔是自然变换. 前函子的复合是由一个余端给出的.

\section{单子}

到现在为止, 你对单子的定义 —— 它是自函子范畴里的一个幺半群 —— 应该已经相当熟悉了. 让我们带着新的
理解重新审视这个定义: 自函子范畴不过是双范畴 $\Cat$ 里自 $1$-胞腔的一个小小的 hom-范畴.
我们知道它是一个幺半范畴: 张量积来自自函子的复合. 幺半群被定义成幺半范畴里的一个物件 —— 这里它将是
一个自函子 $T$ —— 连同两个态射. 自函子之间的态射是自然变换. 一个态射把幺半单位 —— 恒等自函子 —— 映到
$T$:
\[\eta \Colon I \to T\]
第二个态射把 $T \otimes T$ 的张量积映到 $T$. 这个张量积是由自函子的复合给出的, 于是我们得到:
\[\mu \Colon T \circ T \to T\]

\begin{figure}[H]
  \centering
  \includegraphics[width=0.3\textwidth]{images/monad.png}
\end{figure}

\noindent
我们认出来, 这就是定义单子的那两个运算 (在 Haskell 里它们叫 \code{return} 和
\code{join}), 我们也知道幺半群定律在这里变成了单子定律.

现在让我们把自函子这个词从这个定义里完全拿掉. 我们从一个双范畴 $\cat{C}$ 出发, 在其中挑一个
$0$-胞腔$a$. 正如我们早先看到的, hom-范畴 $\cat{C}(a, a)$ 是一个幺半范畴. 因此我们可以
这样定义$\cat{C}(a, a)$ 里的一个幺半群: 挑一个 $1$-胞腔 $T$ 和两个 $2$-胞腔:
\begin{align*}
  \eta & \Colon I \to T         \\
  \mu  & \Colon T \circ T \to T
\end{align*}
它们满足幺半群定律. 我们把\emph{这个}叫作单子.

\begin{figure}[H]
  \centering
  \includegraphics[width=0.3\textwidth]{images/bimonad.png}
\end{figure}

\noindent
这是一个普遍得多的单子定义, 只用到了 $0$-胞腔、$1$-胞腔和 $2$-胞腔. 把它应用到双范畴 $\Cat$ 上, 就
退化成了通常的单子. 但我们来看看它在别的双范畴里会怎样.

我们在 $\cat{Span}$ 里构造一个单子. 我们挑一个 $0$-胞腔, 它是一个集合, 出于马上就会清楚的原因, 我把
它叫作 $Ob$. 接着, 我们挑一个自 $1$-胞腔: 从 $Ob$ 回到 $Ob$ 的一个 span. 它在顶点处有一个集合, 我把
它叫作 $Ar$, 并配有两个函数:
\begin{align*}
  dom & \Colon Ar \to Ob \\
  cod & \Colon Ar \to Ob
\end{align*}

\begin{figure}[H]
  \centering
  \includegraphics[width=0.3\textwidth]{images/spanmonad.png}
\end{figure}

\noindent
我们把集合 $Ar$ 的元素叫作 ``箭头''. 如果我再让你把 $Ob$ 的元素叫作 ``物件'', 你也许就能猜到这是
要通向哪里了. 那两个函数 $dom$ 和 $cod$ 给一个 ``箭头'' 指派定义域和陪域.

要把我们这个 span 变成一个单子, 需要两个 $2$-胞腔: $\eta$ 和 $\mu$. 这里的幺半单位是从 $Ob$ 到
$Ob$、顶点在 $Ob$ 且两个函数都是恒等函数的平凡 span. $2$-胞腔 $\eta$ 是顶点 $Ob$ 与 $Ar$ 之间的
一个函数. 换句话说, $\eta$ 给每个 ``物件'' 指派一个 ``箭头''. $\cat{Span}$ 里的
$2$-胞腔必须满足交换条件 —— 在这里是:
\begin{align*}
  dom & \circ \eta = \id \\
  cod & \circ \eta = \id
\end{align*}

\begin{figure}[H]
  \centering
  \includegraphics[width=0.4\textwidth]{images/spanunit.png}
\end{figure}

\noindent
用分量写出来就是:
\[dom\ (\eta\ ob) = ob = cod\ (\eta\ ob)\]
其中 $ob$ 是 $Ob$ 里的一个 ``物件''. 换句话说, $\eta$ 给每个 ``物件'' 指派一个定义域和陪域都是那个
``物件'' 的 ``箭头''. 我们把这个特殊的 ``箭头'' 叫作 ``恒等箭头''.

第二个 $2$-胞腔 $\mu$ 作用在 span $Ar$ 与自身的复合上. 这个复合被定义成一个拉回, 所以它的元素是
来自$Ar$ 的元素的序对 —— ``箭头'' 的序对 $(a_1, a_2)$. 拉回条件是:
\[cod\ a_1 = dom\ a_2\]
我们说 $a_1$ 与 $a_2$ 是 ``可复合的'', 因为其中一个的定义域是另一个的陪域.

\begin{figure}[H]
  \centering
  \includegraphics[width=0.5\textwidth]{images/spanmul.png}
\end{figure}

\noindent
$2$-胞腔 $\mu$ 是一个函数, 它把一对可复合的箭头 $(a_1, a_2)$ 映成来自 $Ar$ 的单个箭头 $a_3$. 换句话说
,
$\mu$ 定义了箭头的复合.

很容易验证, 单子定律对应于箭头的恒等律和结合律. 我们刚刚定义的恰好是一个范畴 (请注意, 是一个小范畴,
其中物件和箭头构成集合).

所以说, 归根结底, 一个范畴不过是 span 双范畴里的一个单子.

这个结果了不起的地方在于, 它把范畴放到了与单子、幺半群这些代数结构同等的地位上. 身为范畴并没有
什么特别之处. 它不过是两个集合和四个函数. 事实上我们甚至不需要单独拿一个集合来放物件, 因为物件可以
与恒等箭头相认同 (它们一一对应). 所以它其实只是一个集合加上几个函数. 考虑到范畴论在全部数学中扮演的
枢纽角色, 这是一个非常令人谦卑的认识.

\section{挑战}

\begin{enumerate}
  \tightlist
  \item
        对双范畴中定义为自 $1$-胞腔复合的张量积, 推导出单位律和结合律.
  \item
        验证 $\cat{Span}$ 中一个单子的单子定律, 对应于所得范畴中的恒等律和结合律.
  \item
        证明 $\cat{Prof}$ 中的一个单子, 是一个在物件上取恒等的函子.
  \item
        $\cat{Span}$ 中一个单子的单子代数是什么?
\end{enumerate}

\section{参考资料}
\begin{enumerate}
  \tightlist
  \item
        \urlref{https://graphicallinearalgebra.net/2017/04/16/a-monoid-is-a-category-a-category-is-a-monad-a-monad-is-a-monoid/}{Paweł Sobociński 的博客}.
\end{enumerate}
